\paragraph{The same sections over the number field.}
The first term of \eqref{eq:geometric-four-terms} uses a
discriminant at the generic intersection. We identify that
discriminant over the number field, using the same sections
as in the branch computations, before applying the global
height formula.
\hypertarget{source-slice-orders}{}
Put \(F=K(X)\), \(L=\prod_Z L_Z\) and \(\delta=\dim_F L\).

\begin{lemma}[One short basis before localization]
\label{lem:slice-short-basis-before-localization}
With a weight frame \(\tau\) nonzero at every point of the
geometric generic cut,
one can select \(\delta\) short forms from
\(s_F(g_0g_1)^{frw_F.n}\operatorname{Sh}(w_F)\) whose restrictions \(b_i\) give a basis over \(F\),
\[
 b_i=(t'_i/\tau)|_L,\qquad L=\bigoplus_{i=1}^\delta F b_i.
\]

The same list works after extension to every \(K_\lambda(X)\),
with the light columns and \(\eta\) fixed.
\end{lemma}
\begin{proof}
Embed \(K\) in \(\mathbb C\) and specialize \(X\) to algebraically
independent complex numbers over its image. This embeds \(F\) in
\(\mathbb C\). Separability and the chart identify
\(L\otimes_F\mathbb C\) with the product over the \(\delta\) distinct
generic intersection points \(z\). Interpolation of
\hyperref[par:slice-interpolation]{arbitrary values} at weight
\(w_F\), with vanishing low modes, gives the delta functions there.
The bounded rational family spans this section space, so choose
\(\delta\) of its members with \(\det(t'_i(z))_{z,i}\ne0\) in
the chosen nonzero frames. In the \(\tau\)-frame,
\[
 \det(b_i(z))_{z,i}
    =\left(\prod_z\tau(z)^{-1}\right)\det(t'_i(z))_{z,i}\ne0 .
\]
Thus the \(F\)-linear map \(F^\delta\to L\) with columns \(b_i\)
is invertible and remains so over \(K_\lambda(X)\). Clearing
polynomial denominators gives the same assertion for rational
section identities. Only the short restrictions were selected.
\label{cor:slice-orders-1}
\end{proof}

\begin{proposition}[Comparison of the two branch lattices]\label{prop:slice-completed-branch-comparison}
Use the one list of short forms from
Lemma~\ref{lem:slice-short-basis-before-localization}, and suppose
their branch orders are in the pulled-back integral frames.
The integral closures over the two generic coefficient DVRs satisfy
\[
 B'_Z=R_1\otimes_{R_0} B_Z.
\]
 The branches correspond
with the same ramification indices and residue degrees. In
the shifted coordinates of the short forms this gives
the local length inequality and the auxiliary order identity
in \eqref{eq:slice-order-comparison}. The auxiliary identity
is used before the definite last lift.
\end{proposition}

\begin{proof}
\proofstep{Comparing the two integral lattices.}
We compare the two computations of the integral lattice. Put
\[
 R_0=(\mathcal O_{K,\lambda}[X])_{(\pi)},\qquad
 R_1=(\widehat{\mathcal O}_{K,\lambda}[X])_{(\pi)},\qquad
 F_1=\operatorname{Frac}R_1 .
\]
Here \(\pi\) can be taken to generate the maximal ideal of
\(\mathcal O_{K,\lambda}\). Both DVRs have residue field
\(k_\lambda(X)\) and ramification index one over the
coefficient DVR. Let \(B_Z\) be the integral closure of
\(R_0\) in \(L_Z\). It is a finite torsion-free module
over \(R_0\). For example, take an integral field basis.
The trace dual of its span contains every integral element,
by separability and integrality of trace. Each localization
of \(B_Z\) at a maximal ideal is a DVR. Its order \(w\) is
normalized to have image \(\mathbb Z\), with ramification
index \(e_w=w(\pi)\) and residue degree \(f_w\).

The integral closure after base change is exactly

\[
  B'_Z=R_1\otimes_{R_0}B_Z
\]
 in \(F_1\otimes_F L_Z\).
In fact this is finite free over \(R_1\), and its generic
fibre is a product of separable fields. The special fibres identify, because both coefficient rings
have residue \(k_\lambda(X)\):
\[
 (R_1\otimes_{R_0}B_Z)/\pi
       =(R_1/\pi)\otimes_{R_0/\pi}(B_Z/\pi)
       \simeq B_Z/\pi .
\]
This bijects the primes above \(\pi\), even if the generic
fibre splits into several fields. At corresponding
primes localize on both sides; write \(t\) for a uniformizer
of \((B_Z)_w\).

Since \(\pi\) is a unit times \(t^{e_w}\)
there, the quotient of the local ring of \(B'_Z\) by \(t\)
is its residue field. This follows by making that quotient
already in the common localized special fibre. Thus the
maximal ideal is principal. The local ring has dimension
one, since the extension is finite over a DVR and \(\pi\)
is a nonzerodivisor. It is a regular local ring, hence a DVR.
The generic localizations are fields. This proves normality
and the asserted equality with the integral closure, as well
as preservation of both \(e_w\) and \(f_w\).
For corresponding DVRs \(T_w\) and \(T'_{w'}\) on the two models, the branch calculation gives
\[
  e_{w'}=e_w,\qquad f_{w'}=f_w .
\]
 In particular
the branch correspondence counts places even if a field
splits after extending coefficients. \cite{dvr-slice}

\proofstep{Returning to the integer branch order.}
For each branch identified in this way let \(n_w(g)\)
be the integer order of the section \(g\) in an integral
frame of its weight. It is infinite for the zero reading.
To compare with the order on \(L_Z\), restrict the valuation
ring to \(\Omega_Z\). Over this valuation ring the generic
point extends by properness. The induced valuation on
\(\Omega_Z\) is discrete up to a positive integer factor
since \(\pi\) has nonzero order.
Indeed \(v(\Omega_Z^\times)\) is a nonzero subgroup of \(\mathbb Z\) containing \(v(\pi)\). For its positive generator \(b\) this gives
\[
  v(\Omega_Z^\times)=b\mathbb Z,\qquad b\mid v(\pi).
\]

Taking the saturated
flag lattices and a generator of the theta line gives an
integral frame there. Its pullbacks to the two branch
valuations are integral frames. Uniqueness in the valuative
criterion compares the underlying extended points and
the frame changes in the same integral flag have unit
diagonals. Thus for any two nonvanishing sections of the
same weight,
\[
 n_w(g)-n_w(\tau)=w(g/\tau),\qquad
 n_w(g)=e_\zeta\operatorname{val}_{y_\zeta}(g)
\]
at the corresponding branch. This argument fixes the
absolute order of a section as well as the order of a
quotient. It is why the adjusted lattices can be compared
at weights other than the light weight.
\label{cor:slice-orders-2}

\proofstep{Comparing shifted lattices in one trivialization.}
Take \(\tau\) of weight \(w_F\), nonzero on each component,
and choose \(d_Z\in L_Z^\times\) with
\(w(d_Z)=n_w(\tau)\) at all places above \(\lambda\),
by weak approximation. The shifted lattice over \(R_0\)
is \(B(\tau)=\prod_Z d_Z^{-1}B_Z\).
The \(b_i\) belong to it by integrality of the \(t'_i\).
Changing \(\tau\) rescales both the ball and the \(b_i\)
componentwise. Under the coordinate map over \(R_1\)
the local lattice of restrictions with nonnegative point
orders lies in \(R_1\otimes B(\tau)\), and the short
columns map to \(1\otimes b_i\). Hence
\begin{equation}\label{eq:slice-order-comparison}
 I_\lambda\le I^{\mathrm{ff}}_\lambda
 :=\operatorname{length}_{R_0}
       (B(\tau)/\textstyle\sum_i R_0 b_i),\qquad
 N_{\eta,\lambda}=\sum_{Z,w\mid\lambda} f_w n_w(\eta).
\end{equation}
The map on lengths uses flatness and the common residue
field. The second assertion is used only before the final
lift. The integer order already includes the ramification
factor relative to the normalized point valuation.

\end{proof}

\paragraph{A discriminant for the shifted lattice.}

After changing the frame, the discriminant in the index formula is that of the shifted branch lattice. Its change is measured by the determinant of the frame map.

\begin{lemma}[Trace discriminant after the frame shift]\label{lem:shifted-slice-discriminant}
For the branch and shifted lattices of
Proposition~\ref{prop:slice-completed-branch-comparison},
let \(\Delta\) be the discriminant of the trace matrix of the
\(b_i\), and \(t_\lambda\) that of the integral closure in
valuation units. Then the normalized Gauss order satisfies
\[
 v_\lambda(\Delta)+2\sum_{Z,w\mid\lambda}f_w n_w(\tau)
                  =2I_\lambda^{\mathrm{ff}}+t_\lambda .
\]
An empty slice has \(\Delta=1\) and both order contributions
zero.
\end{lemma}

\begin{proof}
Let \(\Delta=\det(\operatorname{Tr}_{L/F}(b_i b_j))\).
The trace pairing of a product of separable fields is
nondegenerate. With an integral basis of \(\prod_Z B_Z\)
let \(t_\lambda\ge0\) be the order of its discriminant.
Write \(C\) for the matrix of \(d b_i\) in an integral
basis of the closure, where \(d=(d_Z)_Z\).
The trace matrix in the new basis is \(C^{\mathsf T}GC\),
with \(G\) the trace matrix of the integral basis. Its
determinant has valuation \(t_\lambda+2I_\lambda^{\mathrm{ff}}\)
by Smith form. The change from \(b_i\) to \(d b_i\)
has determinant \(\mathrm N_{L/F}d\), so the same trace
determinant is \((\mathrm N_{L/F}d)^2\Delta\).

Taking valuations gives
\begin{equation}\label{eq:slice-discriminant-order}
 v_\lambda(\Delta)+2\sum_{Z,w\mid\lambda} f_w n_w(\tau)
     =2I^{\mathrm{ff}}_\lambda+t_\lambda .
\end{equation}
Here the order of the norm is the sum of the orders times
the residue degrees. For an integral \(d\) this follows by
filtering the quotient by \(d\) into the residue fields
of the local DVRs; division gives the formula for
fractional \(d\). The Gauss order of the single rational
function \(\Delta\) thus bounds the first price after the
frame correction. An empty slice has zero lattice, empty
discriminant one and zero contributions in both formulas.
\cite{dvr-slice}\label{cor:slice-orders-3}
\end{proof}
